\documentclass[hidelinks,12pt]{article}
\usepackage[left=0.25cm,top=1cm,right=0.25cm,bottom=1cm]{geometry}
%\usepackage[landscape]{geometry}
\textwidth = 20cm
\hoffset = -1cm
\usepackage[utf8]{inputenc}
\usepackage[spanish,es-tabla, es-lcroman]{babel}
\usepackage[autostyle,spanish=mexican]{csquotes}
\usepackage[tbtags]{amsmath}
\usepackage{nccmath}
\usepackage{amsthm}
\usepackage{amssymb}
\usepackage{mathrsfs}
\usepackage{graphicx}
\usepackage{subfig}
\usepackage{caption}
%\usepackage{subcaption}
\usepackage{standalone}
\graphicspath{{Imagenes/}{../Imagenes/}}
\usepackage[outdir=./Imagenes/]{epstopdf}
\usepackage{siunitx}
\usepackage{physics}
\AtBeginDocument{\RenewCommandCopy\qty\SI}
\ExplSyntaxOn
\msg_redirect_name:nnn { siunitx } { physics-pkg } { none }
\ExplSyntaxOff
\usepackage{color}
\usepackage{float}
\usepackage{hyperref}
\usepackage{multicol}
\usepackage{multirow}
%\usepackage{milista}
\usepackage{anyfontsize}
\usepackage{anysize}
%\usepackage{enumerate}
\usepackage[shortlabels]{enumitem}
\usepackage{capt-of}
\usepackage{bm}
\usepackage{mdframed}
\usepackage{relsize}
\usepackage{placeins}
\usepackage{empheq}
\usepackage{cancel}
\usepackage{pdfpages}
\usepackage{wrapfig}
\usepackage[flushleft]{threeparttable}
\usepackage{makecell}
\usepackage{fancyhdr}
\usepackage{tikz}
\usepackage{bigints}
\usepackage{tcolorbox}
\tcbuselibrary{breakable}
\usepackage{scalerel}
\usepackage{pgfplots}
\usepackage{pdflscape}
\usepackage{enumitem}
\pgfplotsset{compat=1.16}
\spanishdecimal{.}
\renewcommand{\baselinestretch}{1.5}
\def\scaleint#1{\vcenter{\hbox{\scaleto[3ex]{\displaystyle\int}{#1}}}}
\def\scaleoint#1{\vcenter{\hbox{\scaleto[3ex]{\displaystyle\oint}{#1}}}}
\def\scaleiint#1{\vcenter{\hbox{\scaleto[3ex]{\displaystyle\iint}{#1}}}}
\def\scaleiiint#1{\vcenter{\hbox{\scaleto[3ex]{\displaystyle\iiint}{#1}}}}
\def\bs{\mkern-12mu}

\newcommand{\Cancel}[2][black]{{\color{#1}\cancel{\color{black}#2}}}

% \newcommand{\qed}{\tag*{$\blacksquare$}}
\renewcommand{\qed}{\hfill\blacksquare}

\newcommand{\pderivada}[1]{\ensuremath{{#1}^{\prime}}}
\newcommand{\sderivada}[1]{\ensuremath{{#1}^{\prime \prime}}}
\newcommand{\tderivada}[1]{\ensuremath{{#1}^{\prime \prime \prime}}}
\newcommand{\nderivada}[2]{\ensuremath{{#1}^{(#2)}}}

\title{Primer examen parcial - Temas 1 y 2 \\[0.3em]  \large{Matemáticas Avanzadas de la Física}\vspace{-3ex}}
\author{M. en C. Gustavo Contreras Mayén}
\date{ }

\begin{document}
\vspace{-4cm}
\maketitle

\fontsize{14}{14}\selectfont

\textbf{Indicaciones: } Se te pide gentilmente que resuelvas de manera detallada, clara y ordenada los siguientes ejercicios, el puntaje que otorga cada enunciado es de \textbf{1 punto}. En caso de que requieras apoyarte en alguna propiedad, si fue vista en clase, solo indícalo, pero si esa propiedad aunque esté relacionada al ejercicio y no se haya mencionado en clase, habrá que demostrarla debidamente.

\section{Tema 1. La física y la geometría.}
\begin{enumerate}
\item Las ecuaciones de Navier-Stokes para el flujo de un fluido incompresible
    \begin{align*}
    - \curl{( \mathbf{v} \cp (\curl{\vb{v}}))} =  \dfrac{\eta}{\rho_{0}} \, \laplacian{(\curl{\vb{v}})}
    \end{align*}
    Donde $\eta$ es la viscosidad y $\rho_{0}$ la densidad del fluido. Para un flujo axial dentro de un cilindro, consideremos que la velocidad $\mathbf{v}$ está dada por
    \begin{align*}
        \mathbf{v} =  \vu{z} \, v (\rho)
    \end{align*}
    Considera que $\curl{( \mathbf{v} \cp (\curl{\vb{v}}))} = 0$ para este valor de $\mathbf{v}$.
    \\
    Demuestra que
    \begin{align*}
        \laplacian{( \curl{\vb{v}} )} = 0
    \end{align*}
    nos lleva a la ecuación diferencial
    \begin{align*}
        \dfrac{1}{\rho} \, \dv{\rho} \left( \rho \, \dv[2]{v}{\rho} \right) -  \dfrac{1}{\rho^{2}} \, \dv{v}{\rho} = 0
    \end{align*}
    y que la siguiente expresión es solución de la misma ecuación diferencial
    \begin{align*}
        v = v_{0} + a_{2} \, \rho^{2}
    \end{align*}
    \item Escribe los operadores diferenciales: $\grad{\varphi}$, $\div{\vb{A}}$, $\curl{\vb{A}}$, $\laplacian{\varphi}$ en coordenadas esferoidales prolatas, calculando los respectivos factores de escala.
    \item Demuestra que:
    \begin{align*}
        \mathlarger{\Gamma} \left( \dfrac{1}{2} - n \right) \, \left( \dfrac{1}{2} + n \right) = (-1)^{n} \, \pi
    \end{align*}
    donde $n$ es un entero.
    \end{enumerate}

\section{Tema 2. Primeras técnicas de solución.}

\begin{enumerate}
    \item Resuelve con el método de separavión de variable la ecuación de calor:
        \begin{align*}
            u_{t} = 12 \, u_{xx} \hspace{1.5cm} \text{en } 0 < x < \pi, t > 0
        \end{align*}
        sujeto a las siguientes condiciones de frontera e iniciales:
        \begin{align*}
            u_{x} (0, t) &= u_{x} (\pi, t) = 0 \hspace{1.5cm} t \geq 0, \\
            u (x, 0) &= 1 + \sin^{3} x \hspace{1.5cm} 0 \leq x \leq \pi
        \end{align*}
    \item Para el caso especial en donde no hay dependencia en la coordenada azimutal, del estudio del ion molecular del hidrógeno $(H2^{+})$ en mecánica cuántica, se llega a la ecuación
        \begin{align*}
            \dv{\eta} \left[ (1 - \eta^{2} ) \dv{u}{\eta} \right] + \alpha \, u + \beta \, \eta^{2} \, u = 0
        \end{align*}
        Desarrolla una solución en series de potencias para $u(\eta)$. Evalúa los primeros tres coeficientes no nulos en términos de $a_{0}$. Demuestra que la ecuación de índices es $k \, (k - 1) = 0$
        \\
        La expresión a la que debes de llegar es:
        \begin{align*}
            u_{k=1} &=  a_{0} \: \eta \left\lbrace 1 - \dfrac{2- \alpha}{6} \, \eta^{2} + \left[ \dfrac{(2-\alpha)(12-\alpha)}{120} - \dfrac{\beta}{20} \right] \eta^{4} + \ldots \right\rbrace
        \end{align*}
    \item Demuestra que las siguientes ecuaciones diferenciales tienen singularidades en los puntos que se indican:
        \begin{table}[H]
            \centering
            \begin{tabular}{l c c}
                & Sing. regular & Sing. irregular \\
                Hipergeométrica & & \\
                $x (x-1) \sderivada{y} + [(1 + a + b) x - c] \pderivada{y} + a b y = 0$ & $0, 1$ & $--$ \\
                Chevyshev & & \\
                $(1 - x^{2}) \sderivada{y} - x \pderivada{y} + n^{2} y = 0$ & $- 1, 1$ & $--$ \\
                Laguerre & & \\
                $x \sderivada{y} + (1 - x) \pderivada{y} + a y = 0$ & $0$ & $\infty$
            \end{tabular}
        \end{table}
    \item Desarrolla los tres primeros términos no nulos en la serie de potencias en torno a $x = 0$, para la solución general en la siguiente ecuación, con $x > 0$:
        \begin{align*}
            x^{2} \, \sderivada{y} + x \, \pderivada{y} + x^{2} \, y = 0
        \end{align*}
        Recuerda que debes de obtener la primera solución mediante el método de Frobenius y para la segunda solución, considera de la ecuación de índices, si las raíces $r_{1}$ y $r_{2}$ son tales que:
        \begin{itemize}
            \item Si $r_{1} = r_{2}$ entonces la segunda solución linealmente independiente es de la forma:
                \begin{align*}
                    y_{2} (x) = y_{1} \, \ln(x) + \sum_{n=1}^{\infty} b_{n} \, x^{n+r_{1}}
                \end{align*} 
            \item Si $r_{1} - r_{2}$ es un entero positivo, entonces la segunda solución linealmente independiente es de la forma:
            \begin{align*}
                y_{2} (x) = C \, y_{1} \, \ln(x) + \sum_{n=0}^{\infty} b_{n} \, x^{n+r_{2}} \hspace{1cm} b_{0} \neq 0
            \end{align*}
            donde $C$ es una constante que podría anularse.
        \end{itemize}
    \item Una solución para la ecuación de Chebychev
        \begin{align*}
            (1 -x^{2}) \, \sderivada{y} - x \, \pderivada{y} + n^{2} \, y = 0
        \end{align*}
        para $n = 1$ es $y_{1}(x) = x$. Ocupando la doble integral del Wronskiano calcula la segunda solución $y_{2}(x)$.
    \item Para el problema:
        \begin{align*}
            \sderivada{y} - k^{2} \, y = f (x), \hspace{1cm} y (0) = 0, \hspace{0.5cm} \pderivada{y} (0) = 1
        \end{align*}
        \begin{enumerate}[label=\alph*)]
            \item Encuentra la función de Green.
            \item Usa la función de Green para resolver: $\sderivada{y} - y = e^{-x}$
            \item Usa la función de Green para resolver: $\sderivada{y} - 4 \, y = e^{2x}$
        \end{enumerate}
        Revisa que tienes el caso general y luego dos EDO del mismo tipo.
\end{enumerate}

\end{document}